\documentclass[11pt]{article}
% Research investigation; no security theorem for the concrete candidate is claimed.
\usepackage[T1]{fontenc}
\usepackage[utf8]{inputenc}
\usepackage[margin=1in]{geometry}
\usepackage{microtype}
\usepackage{amsmath,amssymb,amsthm,mathtools}
\usepackage{enumitem,booktabs,graphicx,algorithm,algpseudocode,xcolor}
\usepackage[hidelinks]{hyperref}
\usepackage[nameinlink,noabbrev]{cleveref}
\usepackage[backend=biber]{biblatex}
\addbibresource{cryptobib/abbrev3.bib}
\addbibresource{cryptobib/crypto.bib}
\addbibresource{investigation.bib}
\theoremstyle{plain}
\newtheorem{theorem}{Theorem}[section]
\newtheorem{lemma}[theorem]{Lemma}
\newtheorem{proposition}[theorem]{Proposition}
\newtheorem{corollary}[theorem]{Corollary}
\theoremstyle{definition}
\newtheorem{definition}[theorem]{Definition}
\newtheorem{assumption}[theorem]{Assumption}
\theoremstyle{remark}
\newtheorem{remark}[theorem]{Remark}
\newcommand{\secpar}{\lambda}
\newcommand{\bits}{\{0,1\}}
\newcommand{\getsr}{\mathrel{\overset{\$}{\leftarrow}}}
\newcommand{\negl}{\operatorname{negl}}
\newcommand{\poly}{\operatorname{poly}}
\newcommand{\Adv}{\operatorname{Adv}}
\newcommand{\A}{\mathcal A}
\newcommand{\B}{\mathcal B}
\newcommand{\F}{\mathbb F}
\newcommand{\Ftwo}{\mathbb F_2}
\newcommand{\bad}{\perp}
\newcommand{\Enc}{\mathsf{Enc}}
\newcommand{\Dec}{\mathsf{Dec}}
\newcommand{\Gen}{\mathsf{FGen}}
\newcommand{\Punc}{\mathsf{FPunc}}
\newcommand{\Rec}{\mathsf{FRec}}
\newcommand{\Ser}{\mathsf{ser}}
\DeclarePairedDelimiter{\abs}{\lvert}{\rvert}
\DeclarePairedDelimiter{\norm}{\lVert}{\rVert}
\setlist{nosep}
\emergencystretch=2em
\title{Post-Quantum Stacking:\\Candidate Specification and Security Investigation}
\author{}
\date{22 October 2026}
\begin{document}
\maketitle
\begin{abstract}
We specify a delayed-relation candidate commitment using multiple correlated
GGM forests, and identify an index-disclosure mechanism in its unmasked
coefficient response. An explicit four-entry example
distinguishes two compatible indices with constant advantage, without
cryptographic assumptions. Repeated punctured openings also expose complementary
shares. An adaptive pre-tree-key construction also shows that the puncturing
structure alone does not enforce common witness sums across trees.
These issues precede any proof of stacking security. We give typed
algorithms, honest-execution identities, conditional soundness bounds, security
experiments, and unresolved obligations. The source algorithms and revised stacking definitions are inspected directly;
their security theorems do not establish security of this adaptation.
We do not claim a secure equivocal commitment or stacking compiler.
\end{abstract}

\section{Scope, provenance, and stopping points}\label{sec:scope}
The authoritative input is the local research state
\texttt{research-notes/projects/Post-Quantum Stacking.md}, read on
22 October 2026. The original research-state file is unchanged.
The commitment is generated before the opening vector is known, and we retain
the prescribed order
\begin{equation}\label{eq:order}
 c\longrightarrow X\longrightarrow(\boldsymbol\alpha,\boldsymbol\beta)
 \longrightarrow\pi_1\longrightarrow J\longrightarrow\rho.
\end{equation}
All messages after $c$ belong to opening. There is no commitment key with a
preselected binding coordinate in this candidate. Binding is targeted per
commitment, not across all commitments under common parameters.

\paragraph{Sources and provenance.}
The source copies inspected are the 37-page Huth--Joux revision of ePrint
2024/1537 (archive revision dated 5 September 2025), and the 57-page revised
\emph{Stacking Sigmas}, ePrint 2021/422~\cite{investigation:HJ,investigation:Stacking}.
They were retrieved through the ePrint mirror after the main PDF endpoint
returned an access error. We inspect Huth--Joux Sections 2.4--5, Figures
3, 7, and 8, and Appendix A; Stacking Sigmas Definition 4, Figure 5,
Theorem 5 and its proof, and Appendix C; and the relevant definitions of
Avitabile et al.~\cite{investigation:ABFV}.
The formalization below completes serialization and domain labels explicitly.
A previous local manuscript draft in the user's ChatGPT export directory
provided preliminary algebraic examples, independently checked here.
The authoritative research state remains the specified research-notes file.

\paragraph{Status labels.}
\emph{Established algebra} means proved here from stated equations.
\emph{Conditional} identifies a missing forest or distribution hypothesis.
\emph{Unresolved} marks an unproved cryptographic claim. A \emph{proposed change} is not part of the retained candidate.
The most consequential findings are the coefficient leak in
\cref{sec:leak} and the pre-tree counterexample in \cref{sec:pretreeattack}.
No repair is silently substituted for the candidate.

\section{Parameters and types}\label{sec:types}
Let $\lambda\geq1$ be the computational security parameter and the raw
forest-block length, $h\geq1$ the value length, $n\geq2$ the number of entries,
$q=\lceil\log_2 n\rceil\leq\lambda$, $m=\lceil h/\lambda\rceil$, and $t=m+1$.
Indices are zero-based: $[v]=\{0,\ldots,v-1\}$ for positive integers $v$.
Let $D,\tau\geq1$ be tree parameters, $N=2^D$, and $M=\tau D$.
Generation costs include $2^\tau$ pre-tree leaves; an asymptotic parameter
family intended to be efficient must have $2^\tau+\tau2^D$ polynomially
bounded in $\lambda$ and the public input length.
No concrete secure choice of $D,\tau$ is asserted. Requiring the stronger
concrete target $M\geq\lambda$ forces $\max(D,\tau)\geq\sqrt\lambda$
and hence a superpolynomial tree cost as a function of $\lambda$ alone.
Polynomial tree cost still permits $M=\Theta(\log^2\lambda)$ and
negligible $2^{-M}$, but not $2^{-\lambda}$ via this parameterization.
For example, $\tau=16,D=8$ gives $M=128$, with $2^{16}$ pre-tree leaves
per forest; this is a cost illustration, not a security recommendation.
The evaluation field is $\F=\mathbb F_{2^\kappa}$ with a fixed irreducible
polynomial, fixed binary basis, and canonical field-element serialization;
$Q=|\F|=2^\kappa$. The parameters $\kappa$ and $\lambda$ are distinct.
Set $K=t\lambda$. Fix a salt length $s$ and digest length $\ell$; security
losses from these lengths require separate analysis.

\paragraph{Encoding.}
$\Enc(x,i)=\mathbf w\in\Ftwo^K$ splits $x\in\bits^h$ into $m$ consecutive
$\lambda$-bit blocks, zero-pads the last block, and appends a separate
$\lambda$-bit block containing the $q$ low-to-high bits of $i$, followed
by zeros. Let $\mathbf w_a\in\Ftwo^\lambda$ be block $a$.
$\Dec$ checks all padding bits and checks $i<n$; otherwise it returns
$\bad$. The index is never reduced modulo $n$.
Let $z=(z_1,\ldots,z_K)$ denote the corresponding formal variables over $\F$;
$x_u$ and $b_r$ refer to their specified positions.

\paragraph{A necessary type clarification.}
In the bit-polynomial interpretation, an affine evaluation of a block is a
vector in $\F^\lambda$, not one packed element of $\F$:
\begin{equation}\label{eq:typed}
 \mathbf A_a=\sum_{s'=1}^M\alpha_{s'}\boldsymbol\sigma^{(0)}_{a,s'}
 \in\F^\lambda,\qquad
 \mathbf v_a=\mathbf A_a+\Delta\mathbf w_a.
\end{equation}
Here binary shares are embedded coordinatewise in $\F^\lambda$.
This retains $t$ forests; it does not introduce a forest per bit.
The research state's scalar-block notation must be interpreted this way
to evaluate its bit selection polynomials at the advertised degree.
If instead a $\lambda$-bit block is packed into a subfield, bit extraction is
only $\Ftwo$-linear and generally does not commute with multiplication by
$\alpha\in\F$. One would need different polynomial maps and a new degree
analysis. We explicitly adopt the coordinatewise interpretation, as in the vector-valued evaluations in Huth--Joux Section 3.2.

\paragraph{Random functions and serialization.}
Reserve separate domains for tree expansion, leaf hashing, the global digest,
and the two prospective Fiat--Shamir challenges. All inputs encode the
protocol identifier, parameter encoding, salt, forest index, tree index,
node position, and role where applicable. $\Ser$ is injective and
length-delimited. A hash digest has $\ell$ bits. In the classical ROM these
are independent random-function domains; in the QROM the domains permit
superposition queries. Hash collision resistance alone is not a theorem
about correlated-tree consistency or leakage.

\section{Concrete correlated-forest algorithms}\label{sec:forest}
We retain separate forests for the $t$ witness blocks, following the basic
independent-forest idea of Huth--Joux Section 3.3 rather than their
single-forest offset optimization. The tree-generation and puncturing
rules below instantiate Sections 2.4 and 3.1 and Figures 3, 7, and 8.
The extra leaf expansions for their $\mathbf y,z$, offsets, and associated
consistency response are absent: every witness block here is the directly
correlated part of its own forest. Whether those rules force malicious
common-witness consistency is a separate question; see
\cref{sec:pretreeattack} for a counterexample to that inference.

\subsection{Node generation, numbering, and folds}
All nodes are $\lambda$-bit strings. Tree $j=\mathsf{pre}$ has depth $\tau$;
tree $j\in[\tau]$ has depth $D$. A node at level $l$ and position
$p\in[2^l]$ is $T_{a,j}[l,p]$; level zero is unused. Given two level-one
nodes $(r_0,r_1)$, define $\mathsf{Tree}(a,j,L,r_0,r_1)$ by setting
$T[1,0]=r_0,T[1,1]=r_1$ and, for $1\leq l<L$, setting
\begin{align}
 T[l+1,2p]&=H_{\rm node}(\Ser(par,\mathsf{salt},a,j,l,p,T[l,p])),\label{eq:expand}\\
 T[l+1,2p+1]&=T[l,p]\oplus T[l+1,2p].\nonumber
\end{align}
$H_{\rm node}$ has output length $\lambda$. There is no expansion of the
unused root and no check involving a root value. To expand a disclosed
subtree root, apply the same rule at its original absolute labels.
The XOR of all descendants at any level equals the subtree-root value.

For $y=(y_k)_{k\in[2^L]}$ define
\[
 \mathsf{Fold}_{r}^{(b)}(y)=\bigoplus_{k:\operatorname{bit}_r(k)=b}y_k,
 \qquad 0\leq r<L,
\]
where $\operatorname{bit}_0$ is least significant. Thus the most significant
fold ($r=L-1$) equals the two level-one nodes when folding generated leaves.

\subsection{Generation and commitment}
Sample $\mathsf{salt}\getsr\bits^{2\lambda}$ and independent
$R_a\getsr\bits^\lambda$ for $1\leq a\leq t$.
For each block $\mathbf w_a$:
\begin{enumerate}
\item Generate $T_{a,\mathsf{pre}}=\mathsf{Tree}
(a,\mathsf{pre},\tau,R_a,R_a\oplus\mathbf w_a)$ with leaves $p_{a,k}$.
\item For each $j\in[\tau]$ compute
$r_{a,j}^{(b)}=\mathsf{Fold}_{j}^{(b)}(p_a)$ and generate
$T_{a,j}=\mathsf{Tree}(a,j,D,r_{a,j}^{(0)},r_{a,j}^{(1)})$.
Denote its leaves by $u_{a,j,k}$.
\item Compute the binary share vectors and leaf digests
\begin{align}
 \boldsymbol\sigma_{a,j,r}^{(b)}
    &=\mathsf{Fold}_{r}^{(b)}((u_{a,j,k})_k),\label{eq:fold}\\
 C_{a,j,k}&=H_{\rm leaf}(\Ser(par,\mathsf{salt},a,j,k,u_{a,j,k})),\nonumber\\
 \chi_{a,j}&=\bigoplus_k C_{a,j,k},\nonumber\\
 B_{a,j,r}^{(b)}&=\mathsf{Fold}_{r}^{(b)}((C_{a,j,k})_k).\label{eq:hashfold}
\end{align}
$H_{\rm leaf}$ has $2\lambda$-bit output, retaining the explicit leaf
commitment size used in Huth--Joux's basic Figure 3. We use its commitment
output only, rather than its additional $\mathbf y,z$ expansion outputs.
\end{enumerate}
With lexicographic tuple order $(a,j,r,b)$, compute
\begin{align}
 d_{\rm com}&=H_{\rm com}\bigl(\Ser(par,\mathsf{salt},
             (B_{a,j,r}^{(b)})_{a,j,r,b},(\chi_{a,j})_{a,j})\bigr),\nonumber\\
 c&=(\mathsf{salt},(\chi_{a,j})_{a,j},d_{\rm com}).\label{eq:commitment}
\end{align}
Here $H_{\rm com}$ has output length $2\lambda$ (so $\ell=2\lambda$ for this
instantiation). The global hash extends Figure 3's folded-commitment digest
to all forest indices. Its parameter prefix and explicit serialization are
specified domain-separation choices; this is not a verbatim transcription.
Retain $aux=(x,i,\mathbf w,(T_{a,\mathsf{pre}},(T_{a,j})_j)_a,c)$;
all folds can be recomputed. No input or hash query uses $X$.
For honest generation,
\begin{equation}\label{eq:total}
 \bigoplus_k u_{a,j,k}=r_{a,j}^{(0)}\oplus r_{a,j}^{(1)}=\mathbf w_a,
 \qquad
 \boldsymbol\sigma_{a,j,r}^{(0)}\oplus
 \boldsymbol\sigma_{a,j,r}^{(1)}=\mathbf w_a.
\end{equation}
These are honest invariants, not restrictions on a malicious committer.

\subsection{Puncturing}
For a tree of depth $L$ and leaf $k^*$, its punctured key consists of the
$L$ sibling nodes along the path to $k^*$, one at each level $1,\ldots,L$.
If $p_l=\lfloor k^*/2^{L-l}\rfloor$, the level-$l$ sibling is
$T[l,p_l\mathbin{\mathsf{xor}}1]$. Node labels are implicit in $k^*$.

Given $J=(k_j^*)_{j\in[\tau]}\in([2^D])^\tau$, set
\begin{equation}\label{eq:prechallenge}
 \eta_j=\operatorname{bit}_{D-1}(k_j^*),\qquad
 p^*=\sum_{j=0}^{\tau-1}\eta_j2^j.
\end{equation}
For each forest $a$, $\Punc(aux,J)$ sends the full $\tau$-node
pre-tree key punctured at $p^*$, and the second-level sibling nodes at
levels $2,\ldots,D$ for each $j$.
The omitted level-one sibling is recovered from pre-tree folding.
Write the resulting response as
\[
 \rho=(K_{a,\mathsf{pre}},(K_{a,j}^{\rm trunc})_{j\in[\tau]})_{1\leq a\leq t}.
\]
It contains exactly $t(\tau+\tau(D-1))=tM$ node strings, including the
boundary cases $D=1$ and $\tau=1$. No fresh opening randomness is sampled.

\subsection{Reconstruction and rejection conditions}
$\Rec(par,c,J,\rho)$ executes the following deterministic procedure.
\begin{enumerate}
\item Check the canonical encoding of $c,J,\rho$, all lengths and bounds,
and the exact number of nodes; reject malformed or trailing data.
For each $a$, expand the pre-tree sibling subtrees at their original labels.
These partition $[2^\tau]\setminus\{p^*\}$; reject duplicate, missing,
or unauthorized leaves. This is a structural parsing check.
\item For each $j$ compute the known pre-tree half
\[
 \widehat r_{a,j}^{(1-\eta_j)}=
 \bigoplus_{k:\operatorname{bit}_j(k)=1-\eta_j}p_{a,k}.
\]
Use it as the level-one sibling for tree $T_{a,j}$, together with
$K_{a,j}^{\rm trunc}$ at levels $2,\ldots,D$.
Expand these subtrees to reconstruct all leaves except $k_j^*$.
Check the leaf partition as above. There is \emph{no additional root
consistency test}: the missing path is never reconstructed, and no
pre-tree digest appears in $c$.
\item Compute hashes of all known leaves. Recover the missing leaf digest by
\begin{equation}\label{eq:missinghash}
 \widehat C_{a,j,k_j^*}=\chi_{a,j}\oplus
                   \bigoplus_{k\ne k_j^*}C_{a,j,k}.
\end{equation}
Recompute both digest halves $B_{a,j,r}^{(b)}$ for every fold and recompute
\cref{eq:commitment}. Reject if the digest differs from $d_{\rm com}$.
An XOR-restored digest is not a verified preimage of the missing leaf.
\item Set
\begin{equation}\label{eq:knownbit}
 e_{j,r}=1-\operatorname{bit}_r(k_j^*)
\end{equation}
and return
$\boldsymbol\sigma_{a,j,r}^{(e_{j,r})}=\bigoplus_{k:
\operatorname{bit}_r(k)=e_{j,r}}u_{a,j,k}$ for all $a,j,r$.
Never insert a zero string for an unavailable leaf.
\end{enumerate}
This is the concrete reconstruction rule, rather than an idealized
``valid linkage'' predicate. Figure 8 supplies no independent check of
a pre-tree root or its missing-path node relations. That distinction
is decisive for \cref{sec:pretreeattack}.

\paragraph{Sizes and costs.}
Excluding public parameters, $c$ has $2\lambda(t\tau+2)$ bits,
$\rho$ has $t\tau D\lambda=tM\lambda$ bits, and $\pi_1$ has $d\kappa$
bits. The first challenge has $(M+R)\kappa$ bits in the fully independent
sampling specification; the second has $M$ bits. The opening vector itself
has $nh$ bits. A PRG seed compression of the first challenge is a further
computational variant, not part of the independent-challenge analysis.
Generation and full verification involve $O(t(2^\tau+\tau2^D))$ tree
nodes and leaf hashes, in addition to polynomial evaluation. Synchronized
challenges save no punctured nodes across the independently generated forests.

\subsection{What can and cannot be inherited}\label{sec:forestmissing}
The source yields the XOR-preserving expansion, honest fold identities,
the pre-challenge formula, and the compressed key layout above.
Their honest correctness transfers by applying the same calculation
independently to each $a$, with a common $J$ and $\boldsymbol\alpha$.
Huth--Joux Section 3.4's $\mathbf y,z$ consistency test is specific to its
offset extension; it is not a proof of arbitrary malicious consistency
of this multi-forest commitment.
Their Theorem 1 bounds NSBC protocol soundness, not commitment-time
unique-witness binding. Their Theorem 2 and Appendix A concern signature
unforgeability with classical random-oracle query lists; they do not
state the online extraction target of this manuscript or a QROM theorem.
The candidate changes the relation, challenge data, and first response.
Neither theorem is invoked as a security theorem for this candidate.
In particular, their root-counting discussion treats subset sums as distinct
evaluation points; the distribution qualification in \cref{sec:sound} is
needed for a rigorous adaptation.

\section{Delayed relation and batching}\label{sec:relation}
For each $k\in\{0,\ldots,2^q-1\}$ with bits $k_r$, define
\begin{equation}
 \ell_k(b)=\prod_{r=0}^{q-1}
 \begin{cases}b_r&k_r=1,\\1+b_r&k_r=0.\end{cases}
\end{equation}
For Boolean $b$, these are indicator functions. Let $X\in(\bits^h)^n$ be
chosen only after $c$ and fixed before any opening challenge. Set
\begin{equation}\label{eq:selection}
 P_{X,u}(z)=x_u+\sum_{k=0}^{n-1}X[k]_u\ell_k(b),
 \qquad 0\leq u<h.
\end{equation}
On a correctly encoded witness, all these equations vanish precisely when
$x=X[i]$. Their formal degree is at most $\max(1,q)$.

\paragraph{Encoding constraints are necessary for malicious witnesses.}
For a complete relation on $\F^K$, include the following constraints in addition
to \cref{eq:selection}:
\begin{align}
 z_v^2+z_v&=0 &&(1\leq v\leq K),\label{eq:boolean}\\
 z_v&=0 &&\text{for each padding position},\label{eq:padding}\\
 P_{\rm range}(b):=\sum_{k=n}^{2^q-1}\ell_k(b)&=0
 &&\text{if }n<2^q.\label{eq:range}
\end{align}
For Boolean $b$ exactly one selector is one, so \cref{eq:range} excludes
out-of-range encodings, including when $x=0$.
The Boolean constraints increase the degree bound to
$d=\max(2,q)$. They are an explicit proposed completion, not an unnoticed
change to the research state's bound. If a proved forest-extraction property
guarantees a correctly encoded Boolean witness, they may be omitted; then
the retained selection-only checker uses $d_0=\max(1,q)$.
Raw bitstrings alone do not enforce index range or padding against malicious
generation. Our algorithms support both publicly fixed modes:
\emph{selection-only} and \emph{encoding-checked}. The mode is part of $par$.

Enumerate the chosen constraints deterministically as $P_1,\ldots,P_R$.
The verifier samples independently
\begin{equation}\label{eq:challenges}
 \boldsymbol\alpha\getsr\F^M,\qquad
 \boldsymbol\beta\getsr\F^R,
 \qquad P_\beta=\sum_{v=1}^R\beta_vP_v.
\end{equation}
The homogeneous degree bound $d$ is the public bound for the chosen mode,
even when batching cancels top-degree coefficients. Do not choose $d$ from
the resulting batched polynomial. Challenges may contain zero coefficients;
no rank conditioning or resampling is built into the candidate.

\begin{lemma}[Established algebra: fixed-witness batching]\label{lem:batch}
Fix $X$ and $\mathbf w$ before sampling $\boldsymbol\beta$. If at least one
$P_v(\mathbf w)\ne0$, then
$\Pr[P_\beta(\mathbf w)=0]=Q^{-1}$.
\end{lemma}
\begin{proof}
Choose a nonzero constraint value. Conditional on all other batching
coefficients, exactly one choice of its independent uniform coefficient
makes the sum zero.
\end{proof}
Fixing only $c$ is insufficient to invoke this lemma unless a consistency
argument already fixes its witness independently of $\boldsymbol\beta$.

\section{Affine evaluation and polynomial response}\label{sec:poly}
Concatenate the blockwise folds into
$\boldsymbol\sigma_{s'}^{(b)}\in\Ftwo^K$.
The prover computes
\begin{equation}
 \mathbf A=\sum_{s'=1}^M\alpha_{s'}\boldsymbol\sigma_{s'}^{(0)}\in\F^K.
\end{equation}
After receiving $J$ and reconstructing the corresponding shares, the verifier
computes
\begin{equation}\label{eq:affine}
 \Delta=\sum_{j,r}\alpha_{jD+r+1}e_{j,r},\qquad
 \mathbf v=\sum_{j,r}\alpha_{jD+r+1}
                 \boldsymbol\sigma_{jD+r+1}^{(e_{j,r})}.
\end{equation}
For honest folds, coordinatewise in characteristic two,
\begin{equation}\label{eq:affineidentity}
 \mathbf v=\mathbf A+\Delta\mathbf w.
\end{equation}
Synchronization gives one common scalar $\Delta$, not independent evaluation
points for different forests. This is essential to the proposed polynomial check.

Write $P_\beta(z)=\sum_{\nu:|\nu|\leq d}p_\nu z^\nu$.
Its degree-$d$ homogenization is
\begin{equation}\label{eq:hom}
 \widetilde P_\beta(z,Z)=\sum_{\nu:|\nu|\leq d}p_\nu z^\nu Z^{d-|\nu|}.
\end{equation}
Compute as a formal polynomial, not by numerical interpolation on an
undersized field,
\begin{equation}\label{eq:phi}
 \Phi(T)=\widetilde P_\beta(\mathbf A+T\mathbf w,T).
\end{equation}

\begin{lemma}[Established algebra: leading coefficient]\label{lem:leading}
$\deg\Phi\leq d$ and $[T^d]\Phi=P_\beta(\mathbf w)$.
In particular, if the batched constraint vanishes, $\deg\Phi\leq d-1$.
\end{lemma}
\begin{proof}
For each monomial, the coefficient of $T^d$ in
$(\mathbf A+T\mathbf w)^\nu T^{d-|\nu|}$ is $\mathbf w^\nu$.
Sum with coefficients $p_\nu$. This is a formal identity in any field.
\end{proof}
The retained candidate sends the \emph{unmasked} coefficient vector
\begin{equation}\label{eq:response}
 \pi_1=(L_0,\ldots,L_{d-1}),\qquad
 L(T)=\sum_{r=0}^{d-1}L_rT^r=\Phi(T).
\end{equation}
This is defined for honest valid inputs; the honest prover aborts if
$P_\beta(\mathbf w)\ne0$. A malicious prover can send arbitrary coefficients.
The verifier accepts its polynomial check exactly when
\begin{equation}\label{eq:verify}
 \widetilde P_\beta(\mathbf v,\Delta)=L(\Delta).
\end{equation}
The identity in \cref{lem:leading} is not an index-hiding argument; indeed
\cref{sec:leak} shows why the whole coefficient vector is dangerous.

\section{Complete interactive opening and verification algorithms}
The following are complete for the polynomial and transcript layer, using the concrete generation, puncturing, and reconstruction procedures
in \cref{sec:forest}.
All messages are classical, even when the adversary is quantum.
\begin{description}[style=nextline]
\item[$\mathsf{Setup}(1^\lambda,h,n,\kappa,D,\tau,\mathsf{mode})$.]
Check the restrictions in \cref{sec:types}; fix field representation,
constraint enumeration, salt/digest lengths, all domains, and the forest
algorithm identifier/version. Fix $s=\ell=2\lambda$ and the tree rules
of \cref{sec:forest}. Output $par$, or $\bad$.
\item[$\mathsf{Commit}(par,x,i)$.]
Reject incorrect value length or $i\notin[n]$. Form $\mathbf w=\Enc(x,i)$;
sample the common salt and independent forest coins; run the forest
generation procedure and compute \cref{eq:hashfold,eq:commitment}.
Output $(c,aux)$. Do not read $X$ or any challenge.
\item[$\mathsf{OpenStart}(par,c,aux,X)$.]
Check $aux$ belongs to $c$, $X$ contains exactly $n$ binary strings of length
$h$, and $X[i]=x$. Abort otherwise. Fix $X$ and send it or register it as
the public opening input. If an external party selects $X$, this registration
still occurs before the first challenge. The validity test uses no new witness.
\item[$\mathsf{Challenge}_1(par,c,X)$.]
The honest verifier checks public encodings and lengths, enumerates the
constraints for $X$, samples \cref{eq:challenges}, stores the session prefix,
and sends $(\boldsymbol\alpha,\boldsymbol\beta)$.
\item[$\mathsf{Respond}_1(par,aux,X,\boldsymbol\alpha,\boldsymbol\beta)$.]
Reject malformed or wrong-length challenges. Compute $\mathbf A$, $P_\beta$,
and $\Phi$ as in \cref{sec:poly}; abort if its leading coefficient is nonzero.
Send exactly $d$ canonically encoded field elements $\pi_1$ and retain session
state. No polynomial mask is added.
\item[$\mathsf{Challenge}_2(par,c,X,\boldsymbol\alpha,\boldsymbol\beta,\pi_1)$.]
Check the response length and field encodings. Sample independent uniform
$k_j^*\getsr\bits^D$, store $J$, and send it. This sampling occurs strictly
after $\pi_1$ is fixed and is independent of earlier coins.
\item[$\mathsf{Respond}_2(par,aux,J)$.]
Reject malformed $J$; output $\rho=\Punc(aux,J)$, using the same challenge
in every forest. Puncturing is deterministic from fixed forest state and $J$.
\item[$\mathsf{VerifySession}(par,c,X,\boldsymbol\alpha,
\boldsymbol\beta,\pi_1,J,\rho;st_V)$.]
Reject if any public parameter, message, or serialized field element is
invalid, if any message has trailing data, or if the transcript does not
match the verifier's recorded challenge/session state $st_V$.
Run $\Rec$ and reject on any forest failure. Compute \cref{eq:affine},
reconstruct $P_\beta$, and accept iff \cref{eq:verify} holds. Abort/time-out
means rejection. A new opening uses fresh verifier challenges and the same
committed auxiliary state; no reuse-hiding guarantee follows.
\end{description}
\paragraph{Public verification versus authentic challenges.}
One may check the algebra of a recorded interactive transcript, but without
authenticated challenge provenance this is not a sound noninteractive
opening-verification algorithm. A prover could choose its own easy $J$.
The state check above is part of the interactive experiment, not an extra
cryptographic authenticator hidden inside $c$.

\subsection{Honest correctness}
\begin{proposition}[Honest correctness]
Let the forests be honestly generated by \cref{sec:forest}, and let
$\mathbf w=\Enc(x,i)$ with $X[i]=x$.
Then the interactive opening accepts for every challenge value.
\end{proposition}
\begin{proof}
Node expansion preserves XOR at every level, proving \cref{eq:total}.
The pre-tree punctured key gives precisely the omitted level-one siblings;
all other disclosed sibling subtrees partition the authorized leaves.
The restored digest is therefore the honest missing digest and matches $c$.
The selected constraints all vanish on $\mathbf w$, in either mode.
Thus $P_\beta(\mathbf w)=0$ for every $\boldsymbol\beta$ and
$\pi_1$ is well-defined by \cref{lem:leading}.
Honest reconstruction passes the digest checks and gives
\cref{eq:affineidentity}. Substitution at $T=\Delta$ proves
\cref{eq:verify}, including $\Delta=0$. The latter case is correct but can
be especially easy to exploit in a malicious transcript.
\end{proof}

\subsection{Prospective Fiat--Shamir variant: distinct and unproved}
Let $ctx$ bind the application and its public statement. An explicit candidate
derivation would be
\begin{align}
 (\boldsymbol\alpha,\boldsymbol\beta)
 &=\mathsf{Sample}_1(H_{\rm FS1}(\Ser(par,ctx,c,X))),\nonumber\\
 J&=\mathsf{Sample}_2(H_{\rm FS2}(\Ser(par,ctx,c,X,
              \boldsymbol\alpha,\boldsymbol\beta,\pi_1))).\label{eq:fs}
\end{align}
The sampling functions use a domain-separated expandable oracle stream and
canonical uniform field/bit sampling. The opening would contain $(\pi_1,\rho)$;
verification recomputes both challenges. This is not the interactive protocol:
the prover can grind $c$, $X$, and $\pi_1$, and requery the oracle.
Interactive root bounds do not automatically prove ROM Fiat--Shamir security,
and classical ROM security would not automatically prove QROM security.
Most importantly, Fiat--Shamir cannot remove the coefficient leakage below.

\section{Attack-oriented index-hiding analysis}\label{sec:leak}
\subsection{Candidate-witness testing from one transcript}
For each public index $i$, define $\mathbf w^{(i)}=\Enc(X[i],i)$.
A verifier knows $\Delta$ and $\mathbf v$ from its reconstructed view and can
compute, without any forest secret,
\begin{equation}\label{eq:test}
 F_i(T)=\widetilde P_\beta\bigl(\mathbf v+(T+\Delta)\mathbf w^{(i)},T\bigr).
\end{equation}
For the actual honest index $i^*$, \cref{eq:affineidentity} implies
$F_{i^*}(T)=L(T)$ identically. Every candidate that fails this coefficient
comparison is ruled out. This is an exact transcript test, even if the forest
commitment alone is hiding and even if its punctured view is perfectly simulated.
It is efficient in $n$ and the cost of polynomial computation; the opening
already supplies $n$ entries. It does not require discrete logarithms, hash
inversion, rewinding, or quantum queries.

\subsection{Explicit four-entry distinguishing example}\label{sec:four}
Work in selection-only mode with $n=4$, $h=1$, $q=2$, and $d=2$.
Let $i=b_0+2b_1$ and
\begin{equation}
 X=(0,0,0,1),\qquad P(z)=x+b_0b_1,
 \qquad\widetilde P(z,Z)=xZ+b_0b_1.
\end{equation}
Compare indices $0$ and $1$, both committing to the value $x=0$.
Write the three relevant mask coordinates as $(A_x,A_0,A_1)$ and verifier
coordinates as $(v_x,v_0,v_1)$. For a nonzero batching scalar $\beta$, divide
the sent coefficients by $\beta$. Expanding gives
\begin{equation}\label{eq:quad}
 \Phi(T)=A_0A_1+(A_x+A_0b_1+A_1b_0)T
                 +(x+b_0b_1)T^2.
\end{equation}
For the two candidate indices, $x=b_1=0$, and therefore
\begin{equation}\label{eq:leak}
 \beta^{-1}L_1=
 \begin{cases}
 v_x,&i=0,\\
 v_x+v_1,&i=1.
 \end{cases}
\end{equation}
When $v_1\ne0$, these are distinct. A verifier seeing $\beta\ne0$ and
$v_1\ne0$ recovers the selected index with certainty, using only a single
honest opening. Otherwise it can guess a fair bit. The event and test are public.

\paragraph{Concrete constant-advantage attack.}
The nonzero-mask event has constant probability under the actual honest
forest generation, without a pseudorandomness assumption. Let $a$ be the
index forest and $v_1$ its coordinate for $b_1$. The most significant
pre-tree fold is $(R_a,R_a\oplus\mathbf w_a)$. Thus the second-level tree
$j=\tau-1$ has left root $R_a$, and its most significant fold has
$\boldsymbol\sigma_{a,\tau-1,D-1}^{(0)}=R_a$.
With probability $1/2$, the $b_1$ coordinate of this share is one.
Consequently the vector of $M$ binary zero-half shares for that coordinate
is nonzero with probability at least $1/2$.
Conditional on any nonzero such vector, uniform independent
$\boldsymbol\alpha$ makes $A_1=v_1$ uniform over $\F$; here $b_1=0$
in both challenge worlds. Independently $\Pr[\beta\ne0]=1-Q^{-1}$.
A distinguisher applies \cref{eq:leak} on the separating event and guesses
fairly otherwise. In the standard convention $\Pr[\text{correct}]-1/2$,
its advantage is at least
\begin{equation}\label{eq:advleak}
 \frac14(1-Q^{-1})^2.
\end{equation}
The two challenge pairs commit to the same $x=0$, so this directly
violates static single-opening index hiding, even against an honest verifier.
The attack also applies in the classical ROM and the QROM.

\paragraph{The optional encoding constraints do not repair this example.}
If they are added with independent random batching coefficients, their
homogenized Boolean contribution is $z_v^2+z_vZ$.
On $\mathbf A+T\mathbf w$ with Boolean $\mathbf w$, its coefficient of $T$
is $A_v$. Padding constraints, homogenized to degree two, also contribute
$A_v$ to that coefficient. After subtracting the verifier-computable terms
involving $v_v$, the coefficient has the form
\begin{equation}
 S=b_0(\beta v_1+\gamma_0\Delta),\label{eq:checkedleak}
\end{equation}
where $\gamma_0$ is the batching coefficient of the Boolean constraint for
$b_0$ (there is no range constraint for $n=4$). Terms for other coordinates
cancel or have zero witness coordinate in these two worlds.
Explicitly, $S$ is $L_1$ plus $\beta v_x$ plus the sums of the Boolean
and padding batching coefficients times their corresponding visible
coordinates. The two indices remain distinguishable whenever the displayed multiplier
is nonzero. Conditional on $(v_1,\Delta)\ne(0,0)$, its independent
uniform coefficients make this multiplier nonzero with probability
$1-Q^{-1}$. The same mask argument proves at least the bound in
\cref{eq:advleak} for this mode as well.
Adding encoding checks addresses relation validity, not masking
of the polynomial response.

\subsection{Why a simulator without the index is not yet defined}
A simulator given only $X$ and the verifier's reconstructed affine view must
produce $L$ whose joint distribution is compatible with both challenge
witnesses. Equation~\eqref{eq:leak} makes the two supports disjoint on the
nondegenerate event. Producing one accepting scalar equality at $\Delta$
does not simulate the distribution of all coefficients. This distinguishes
correctness/soundness checks from zero knowledge or index hiding.

\subsection{Repeated puncturing leaks the committed witness}
Even if the coefficient issue were repaired, fresh openings reuse the same
forests. Suppose two honest challenges differ in a bit $e_{j,r}$. The verifier
then obtains both halves of that fold, and for every forest $a$ recovers
\begin{equation}
 \mathbf w_a=\boldsymbol\sigma_{a,j,r}^{(0)}
                       \oplus\boldsymbol\sigma_{a,j,r}^{(1)}.
\end{equation}
Equivalently, two different punctures of one tree reveal all its leaves.
With independent uniform $J,J'$, there is a changed challenge bit with
probability $1-2^{-M}$. Thus, for the concrete honest reconstruction procedure, two complete openings reveal all blocks of $\Enc(x,i)$ with that
probability. Multi-opening \emph{binding} does not itself ask for multi-opening
\emph{hiding}; applications must decide which they require.
This extraction from two responses also does not prove online extraction
from queries available before $c$.

\section{Conditional polynomial soundness}\label{sec:sound}
The useful algebraic bound is conditional on an adversarial commitment
already determining common folds and a witness $\mathbf w$ before
$\boldsymbol\beta$, with \cref{eq:affineidentity} valid for all accepted
challenges. This is a substantive missing forest-consistency hypothesis.

\subsection{The evaluation point is not automatically uniform in the field}
Express the $\alpha_{s'}$ in the chosen binary field basis. For fixed
$\boldsymbol\alpha$, let
\begin{equation}
 r(\boldsymbol\alpha)=\operatorname{rank}_{\Ftwo}
       (\alpha_1,\ldots,\alpha_M),\quad
 V_\alpha=\operatorname{span}_{\Ftwo}\{\alpha_1,\ldots,\alpha_M\}.
\end{equation}
Because independent uniform punctures make $e\in\bits^M$ uniform,
$\Delta$ is uniform on $V_\alpha$, with point probability $2^{-r(\alpha)}$.
It is uniform on $\F$ only if $r(\alpha)=\kappa$.
For $M<\kappa$ this cannot occur. For $M\geq\kappa$,
\begin{equation}
 \Pr[r(\alpha)<\kappa]\leq(2^\kappa-1)2^{-M}
\end{equation}
by a union bound over nonzero dual linear functionals.

\begin{proposition}[Conditional single-session bound]\label{prop:sound}
Assume the fixed-witness/fixed-fold hypothesis above; let the chosen
constraints fail on $\mathbf w$. For honest independent challenges and any
degree-at-most-$(d-1)$ response fixed before $J$, acceptance probability is
at most
\begin{equation}\label{eq:bound}
 \min\left\{1,\; Q^{-1}
      +d\bigl(2^{-M}+(1-2^{-M})Q^{-1}\bigr)\right\}.
\end{equation}
For a fixed first-challenge prefix outside the batching-zero event, its
second-challenge acceptance probability is at most
$\min\{1,d\,2^{-r(\alpha)}\}$.
\end{proposition}
\begin{proof}
Except with probability $Q^{-1}$, \cref{lem:batch} gives
$P_\beta(\mathbf w)\ne0$. The residual
$R(T)=\widetilde P_\beta(\mathbf A+T\mathbf w,T)-L(T)$
then has nonzero leading coefficient and degree $d$, even if $\mathbf A$
and $L$ depend on both first challenges. It has at most $d$ field roots.
Conditional uniformity on $V_\alpha$ proves the fixed-prefix bound.
For independent uniform $\alpha_{s'}$,
\begin{align}
 \mathbb E_\alpha[2^{-r(\alpha)}]
 &=\Pr_{\alpha,e}\left[\sum_{s'}\alpha_{s'}e_{s'}=0\right]\nonumber\\
 &=2^{-M}+(1-2^{-M})Q^{-1}.
\end{align}
The zero bit vector always yields zero; for any nonzero vector the sum is
uniform over $\F$. Add the batching-zero probability. No independence of
the adversary's residual polynomial from $\alpha$ is needed.
\end{proof}
This bound becomes useless when the degree is comparable to the span size.
An invalid opening can have a negligible but nonzero acceptance probability;
extraction must permit this soundness error rather than promise literal
deterministic uniqueness for every possible accepting transcript.

\subsection{Challenge dependencies and failure modes}
If $X$ is selected after $\boldsymbol\beta$, the fixed-constraint batching
argument does not apply. If $\pi_1$ is selected after $J$, a constant
$L(T)=\widetilde P_\beta(\mathbf v,\Delta)$ always satisfies the final equality.
If reconstructed $\mathbf v$ does not come from one committed affine line,
root counting on that line proves nothing about malicious forests.
If the prover selectively aborts depending on $J$, the bound applies to
unconditional session success, not to success conditioned on completion.
If $S$ sequential sessions each satisfy the same conditional bound, a union
bound gives at most $S$ times that error for any invalid accepted opening;
this does not justify concurrent sessions or Fiat--Shamir grinding without
an additional argument.

\section{Forest consistency, binding, and online extraction}\label{sec:binding}
\subsection{Adaptive pre-tree keys do not force a common sum}\label{sec:pretreeattack}
\begin{proposition}[Counterexample to pre-tree consistency]\label{prop:pretree}
For any collection of honestly expanded second-level trees, whose root
pairs and total XORs may be unrelated, a committer can compute $c$ in
\cref{eq:commitment} before $J$ and, for every $J$, produce a response
passing $\Rec$. It need not know one pre-tree linking these trees.
The statement concerns forest acceptance, not the final polynomial test.
\end{proposition}
\begin{proof}
Fix one forest $a$ and arbitrary root pairs $(r_j^{(0)},r_j^{(1)})$.
Expand each second-level tree honestly, compute its hashes and all folds,
and commit to these before $J$. After seeing $J$, its truncated tree keys
are fixed. We must produce a pre-tree punctured key whose known fold $j$
is the target $t_j=r_j^{(1-\eta_j)}$.

For $\tau=2$, let $V=t_1$ be the disclosed opposite level-one node of
the pre-tree. Write $\mathsf{child}_b(V)$ for its $b$ child under
\cref{eq:expand} at that absolute position. Set the disclosed sibling
leaf on the punctured side to
\[
 U=t_0\oplus\mathsf{child}_{1-\eta_0}(V).
\]
The known most significant fold is $V=t_1$, and the known least
significant fold is $U\oplus\mathsf{child}_{1-\eta_0}(V)=t_0$.
The punctured key $(V,U)$ therefore supplies precisely the two desired
second-level root siblings, for all four $(\eta_0,\eta_1)$.
No missing-path parent is disclosed or checked.

For arbitrary $\tau$, construct the punctured-key siblings from top to
bottom. At level $l$ let $r=\tau-l$, and let $E_{l-1}$ be the set of
pre-tree leaves already reconstructed from earlier siblings.
Set the new sibling value to
\begin{equation}\label{eq:inversefold}
 K_l=t_r\oplus
 \bigoplus_{k\in E_{l-1}:\operatorname{bit}_r(k)=1-\eta_r}p_k.
\end{equation}
Expand it at its absolute position. All its descendants have bit $r$
equal to $1-\eta_r$, and their XOR is $K_l$. All later siblings lie
on the punctured side of this bit and cannot change the known fold.
Thus the newly completed fold is $t_r$. Induction yields all $\tau$
target folds. This costs at most the honest pre-tree expansion cost.

Use these pre-tree keys and the actual truncated tree keys. Reconstruction
now returns the originally committed second-level leaves except the
punctures. Every restored hash and global digest is correct. Independent
application to every $a$ proves the claim. No oracle collision, inversion,
or programming is used.
\end{proof}

\paragraph{Consequence and scope.}
For example, choose two second-level trees with respective sums $0^\lambda$
and $10^{\lambda-1}$. Both are supported under one forest commitment for
every challenge. Their binary shares cannot all sum to one witness block.
Therefore the pre-tree does not certify the malicious-prover invariant
\cref{eq:total}, despite the honest-generation linkage.
This is an explicit counterexample to the research state's proposed
structural consistency argument, using the actual basic puncturing layout.
It does \emph{not} yet produce two polynomial-accepted incompatible opening
vectors, nor refute every possible online extractor. A proof might instead
extract several per-tree witnesses and analyze the polynomial test jointly.
It cannot simply assume one common affine line because the keys are
called correlated-forest keys.

\paragraph{An acceptance-probability illustration.}
Take $n=2$, $h=1$, selection-only mode ($d=1$), and one value forest
with one tree having value bit $1$ and all remaining trees value bit $0$.
All index-tree sums encode $0$; all unused bits are zero.
For $X=(0,0)$ the polynomial is $P_\beta=\beta x$.
Send $L_0=\beta A_x$ before $J$. The verifier's residual is
$\beta\Delta_{j_*}$, where
$\Delta_{j_*}=\sum_r\alpha_{j_*,r}e_{j_*,r}$ belongs to the exceptional
tree. The forged pre-tree keys pass forest verification for every $J$.
For $\beta\ne0$, averaged over the honest first challenge, the polynomial
check accepts with probability
\[
 2^{-D}+(1-2^{-D})Q^{-1},
\]
rather than a bound based on $2^{-M}$.
A common witness $(0,0)$ is nevertheless compatible with this particular
opening: this example refutes a claimed all-tree common-line proof, not
the per-commitment binding game. It also shows why a commitment-time
extractor cannot select an arbitrary exceptional tree and claim the
stronger $M$-bit error bound for that choice.

\subsection{A binding attack on the unchecked non-power-of-two extension}
Take $n=3$, $h=1$, and $q=2$. In selection-only mode, a malicious prover
can generate otherwise honest forests for the raw encoded bits $x=0$,
$b_0=b_1=1$, corresponding to the excluded index $3$.
All three in-range selectors vanish, so
\begin{equation}
 P_{X,0}(\mathbf w)=0\qquad\text{for every }X\in\bits^3.
\end{equation}
Consequently, conditional only on honest forest correctness for arbitrary
raw block strings, it can answer both openings $X^{(0)}=(0,0,0)$ and
$X^{(1)}=(1,1,1)$ correctly for every challenge. Their coordinate--value
sets have empty intersection. Rejection of the excluded index in the honest
$\mathsf{Commit}$ API does not restrict a malicious committer.
This breaks per-commitment binding in that unchecked extension. It does
\emph{not} attack the explicitly power-of-two selection relation in the
research state. The encoding-checked mode's range constraint is a proposed
correction for the extension, and must be proved sound against malicious
forests rather than merely appended to an honest-execution argument.

\subsection{Digest binding is a different property from witness consistency}
The global digest authenticates folded hash data against alternative
reconstructions, subject to a suitable random-oracle argument. Distinct
leaf arrays can have identical folded digest XORs as linear data; proving
that efficiently obtainable alternative leaf preimages cannot exploit
this kernel requires more than the collision resistance of the outer hash. It does not,
by itself, show that a maliciously generated collection has
$\bigoplus_k u_{a,j,k}$ independent of $j$, or that it comes from a valid
correlated pre-tree. Hashing arbitrary data cannot establish the algebraic
relation that honest generation was meant to impose.

\paragraph{Counterexample to omitting pre-tree linkage.}
If a verifier merely accepts independent punctured second-level trees, a
malicious committer can generate each from an arbitrary raw block sum
$w_{a,j}$, with different sums in different trees, compute the same digest
format honestly for those leaves, and open each tree consistently in isolation.
Then the reconstructed evaluations satisfy
\begin{equation}
 \mathbf v_a=\mathbf A_a+
 \sum_j\left(\sum_r\alpha_{jD+r+1}e_{j,r}\right)w_{a,j},
\end{equation}
which need not equal $\mathbf A_a+\Delta w_a$ for any common $w_a$.
The explicit construction in \cref{prop:pretree} shows that, for forest
acceptance, the actual pre-tree layout can also be satisfied by this
independent-tree strategy. Additional polynomial or source-protocol
checks may constrain it; their effect needs its own proof.

\subsection{Interactive security experiments}
\begin{definition}[Per-commitment multi-opening binding target]
An efficient adversary outputs one $c$, then conducts polynomially many
opening sessions with an honest verifier, choosing each $X^{(s)}$ after $c$
and before that session's fresh first challenge. For the set $\mathcal S$ of
accepted sessions define
\begin{equation}
 B_s=\{(i,X^{(s)}[i]):i\in[n]\}.
\end{equation}
The adversary wins if $\mathcal S\ne\varnothing$ and
$\bigcap_{s\in\mathcal S}B_s=\varnothing$. Require negligible win probability.
The experiment specifies sequential sessions; a concurrent version is a
separate target. One accepted session alone cannot win.
\end{definition}
No fixed index is required across distinct commitments. This definition
does not by itself provide an efficient extractor, nor a witness selected
at commitment time independently of future challenges.
Pairwise intersections are insufficient: for $n=3$, the vectors
$(0,0,0)$, $(0,1,1)$, and $(1,0,1)$ have pairwise common coordinate--value
pairs, but the intersection of all three corresponding sets is empty.

\begin{definition}[Commitment-time classical online extraction target]
An extractor gets $par,c$ and the classical oracle query/answer transcript
$Q_{\leq c}$ at the moment $c$ is published, without rewinding the committer.
It outputs $w^*\in\bits^h\times[n]$ or $\bad$ before $X$ is selected.
In any polynomial subsequent sequential interaction with the honest verifier,
require negligible probability of an accepted opening $X$ for which
$w^*=\bad$ or $X[i^*]\ne x^*$.
\end{definition}
This permits nonextractable unopenable commitments. It does not require
rejecting every arbitrary bitstring as a commitment at publication time.
It is stronger than the preceding intersection game; a union bound over
sessions gives the implication if extraction failure has the stated bound.

\subsection{A classical extraction program, and its unproved steps}
A natural program searches pre-commitment $H_{\rm com}$ queries for the
serialized input matching $d_{\rm com}$, and associates leaf-hash and tree
queries by the salt and all labels. If it obtains a coherent complete forest,
it computes block XORs, checks common sums across trees, and applies $\Dec$.
It must additionally prove all of the following:
\begin{enumerate}
\item a successfully openable $c$ has an appropriate digest query before publication,
except with a quantified guessing probability;
\item hash-fold inputs are linked to sufficiently many actual leaf preimages,
not merely arbitrary XOR aggregates;
\item some mechanism beyond the current pre-tree keys yields an appropriate
commitment-time witness, or a list of per-tree candidates can be converted
to one witness with a valid continuation bound; \cref{prop:pretree} rules
out the simple common-block premise;
\item a purported index block is valid; otherwise range/padding failures cannot
be ignored because honest commitments happen to encode valid indices;
\item the total extraction and search cost is polynomial in the query count and
parameters, without relying on enumerating exponentially many possible trees.
\end{enumerate}
Querying the digest on arbitrary aggregates is easy; hence finding its
preimage alone does not discharge the second or third items. Two-opening
special extraction is not a substitute for the stipulated pre-$X$ extractor.
All these extraction claims remain unresolved for the concrete candidate.

\subsection{QROM distinction}
A QROM adversary does not expose a classical list of queried inputs;
measuring every query changes its computation. Thus the preceding transcript
search is not a QROM extractor. Don, Fehr, Majenz, and Schaffner
provide an extractable random-oracle simulator (Theorem 4.3) and early
extraction results (Corollaries 4.7--4.8)~\cite{investigation:DFMS}.
Their setting associates a public classical value $t$ with a relation
$f(z,H(z))=t$, and quantifies both extraction and disturbance through
relation-dependent parameters $\Gamma,\Gamma'$ and oracle query counts.
It permits extracting immediately after $t$ is announced; it does not
license reading quantum queries as a classical transcript.

For the outer digest, a natural attempt is $f(z,h)=h$ and
$t=d_{\rm com}$. Even a successful extraction here returns folded hash
strings, not their leaf seeds, a correlated pre-tree, or the witness.
A further extraction composition must link these XOR strings to oracle
preimages and valid per-tree sums, account for its disturbance and cost,
and then justify a single committed witness. The adaptive-key counterexample
survives perfect extraction of all second-level trees, so this is also an
algebraic gap, not just a QROM query-access difficulty.
Huth--Joux Appendix A explicitly maintains classical query and answer
lists and candidate witnesses. Its candidate search occurs during the
forgery experiment, with final offsets available; it neither proves the
pre-$X$ unique-witness claim here nor supplies a quantum-query analysis.
The classical leakage in \cref{sec:leak} also applies against quantum
adversaries; a QROM theorem cannot repair a classical information leak.

\section{Index hiding and composition targets}\label{sec:composition}
\paragraph{A single-opening hiding game.}
The distinguisher chooses two pairs $(x_0,i_0),(x_1,i_1)$ before commitment.
A hidden bit $b$ selects the honest commitment to $(x_b,i_b)$.
After seeing $c$, the distinguisher chooses $X$ such that
$X[i_0]=x_0$ and $X[i_1]=x_1$; failure of this common compatibility condition
ends the experiment with an independent fair guess. It then receives an
honest opening interaction, including the verifier's reconstructed view,
and guesses $b$. A minimal target restricts it to honest challenge sampling;
the malicious-verifier target lets it choose well-formed challenges.
Require negligible advantage. Restricting to pairs $(X[i_0],i_0)$,
$(X[i_1],i_1)$ and fixing $X$ in advance gives a weaker static game.
The example in \cref{sec:four} already separates that static game under its
explicit nondegeneracy condition. Commitment-only hiding is insufficient.
Selective abort behavior must not reveal the selected pair.

\paragraph{Commitment-only hiding is a separate target.}
Require indistinguishability of the marginal commitments generated from any
two legal pairs $(x_0,i_0),(x_1,i_1)$. The displayed coefficient attack uses
the opening and does not by itself distinguish these marginals. A proof
of their hiding would need the correlated-tree and leaf-hash distributions;
none is supplied here. Conversely, the proposed full-opening hiding game
would imply this marginal target for $n\geq2$: for distinct indices choose
a vector compatible with both pairs and discard the opening. For equal
indices and different values, use a hybrid through a third commitment
pair at another index, with one compatible vector for each transition.
This implication is a property of the security definitions, not a hiding
proof for the insecure candidate.

\paragraph{What per-commitment binding can support.}
If an eventual compiler's extractor obtains enough accepted openings of one
\emph{same} $c$, the intersection property supplies at least one common
coordinate and value. If $n$ and the number of openings are polynomial, it
can compute that intersection from the revealed vectors.
This still requires a reduction showing that the openings are obtainable
with appropriate challenge probabilities and correspond to the base
protocol's special-soundness transcripts. An existential intersection gives
neither a pre-challenge index nor a general quantum extractor.

\paragraph{Why cross-commitment extraction is a distinct obligation.}
Suppose a compiler uses $c_1,c_2$ and needs a common binding coordinate.
Per-commitment binding permits $c_1$ to bind at index zero and $c_2$ to bind
at index one, with no reason for them to share an index. This logical
countermodel refutes an implication to cross-commitment consistency.
It is not an attack on a fully specified compiler, which is not yet present.
The distinction can be checked directly against revised Stacking Sigmas
Definition 4 (pp.~12--13). For an adversarially chosen key $ck$, it forbids
a polynomial set $\mathcal T$ of tuples $(v,v',r,r')$ for which
\begin{equation}\label{eq:revisedbinding}
 \left|\bigcup_{(v,v',r,r')\in\mathcal T}
       \{j:v_j\ne v'_j\}\right|>\ell-t,
 \quad
 \mathsf{BindCom}(pp,ck,v;r)=\mathsf{BindCom}(pp,ck,v';r')
\end{equation}
for every tuple. Different tuples may correspond to different commitment
values. Thus at least $t$ coordinates remain unchanged across \emph{all}
listed pairs under that key. The adjacent prose contains an apparent
$t-\ell$ sign typo; the displayed winning condition uses $\ell-t$.
Avitabile et al. Definition 6's fixed-equivocation game explicitly combines
two commitments under the same parameter pair, changing opposite positions
in the two pairs~\cite{investigation:ABFV}.

The revised hiding condition also equates a key/commitment/equivocated-opening
distribution with direct binding commitment under another binding set and
uniform opening randomness. It is stronger than a commitment-only index
hiding condition. A new interactive primitive would need an appropriate
transcript-level analogue, in addition to repairing the present separator.

There is a useful limited positive observation: Theorem 5's one-level
self-stacking special-soundness proof (p.~27) compares two openings of the
\emph{same} first-message commitment. A per-commitment guarantee can provide
a common coordinate in that particular step, if openings are genuinely
sound and obtainable. The recursive stackability argument also needs its
extended honest-verifier simulation, recyclable third-message distribution,
and key consistency interfaces. Those requirements do not follow from this
single extraction step. The threshold setting of Avitabile et al. needs
linked binding positions across commitments; the countermodel above explains
why its fixed-equivocation property cannot be supplied by the present target.
No whole-compiler security theorem is inherited merely because the local
two-opening argument has a common coordinate.

\paragraph{Round structure matters.}
The commitment opening has two verifier challenges and is not itself a
three-move Sigma protocol. A stacking compiler must specify how
\cref{eq:order} interacts with base-protocol challenges, polynomial responses,
repeated openings, and any rewinding/quantum extraction. Declaring the
commitment an equivocal primitive does not remove these messages.
Even a successful binding proof would not establish the necessary hiding
and composition properties.

\section{A proposed masking interface, not an implemented repair}\label{sec:mask}
QuickSilver Section 5, Figure 6, uses the same degree-$d$ leading-coefficient
identity but additionally obtains a VOPE mask from its preprocessing
functionality~\cite{CCS:YSWW21}. Its prover adds correlated random coefficients
to the batched response, and the verifier adds the corresponding evaluation.
Its hybrid-model security theorem does not instantiate that preprocessing
from the present forests or prove security under our challenge order.
This comparison identifies a substantive missing ingredient rather than a
new degree-reduction technique.

Concretely, a candidate correction would require a commitment-time auxiliary
correlation with the following interface. The prover obtains a degree-at-most-
$(d-1)$ polynomial $R(T)$ independent of the witness. At puncturing time the
verifier obtains $B=R(\Delta)$, authenticated against that same earlier
correlation. The prover would send $U(T)=\Phi(T)+R(T)$, and verification would
check
\[
 \widetilde P_\beta(\mathbf v,\Delta)+B=U(\Delta).
\]
For privacy, conditional on the entire verifier view and $B$, the mask
coefficients should have the uniform distribution over the affine space
$\{R:R(\Delta)=B,\deg R<d\}$, up to a proved statistical or computational
error. Then all candidate witnesses compatible with the same evaluation
produce the same conditional distribution of $U$: addition by their
$\Phi$ is a bijection between the corresponding affine spaces.
This is a conditional algebraic simulation argument, not a simulation of
the full forests or an implemented preprocessing construction.

The correlation must be fixed before $X$ and before the polynomial response;
$R(\Delta)$ cannot be supplied as an unauthenticated scalar after $J$, since
that would let a prover make the final equality hold arbitrarily. Hiding of
punctured keys, malicious correlation consistency, extraction, degree-dependent
costs, and possible multi-opening leakage would all require new analysis.
The retained candidate still sends unmasked $L$ and is insecure.

\section{Proposed corrections and unresolved work}\label{sec:next}
The following are candidates for discussion, not silently applied repairs.
\begin{enumerate}
\item Address the coefficient leak before proving index hiding. A possible
direction is auxiliary, witness-independent masking material committed
before $X$, enabling a simulator to hide the lower coefficients while
preserving the final evaluation. Simply adding an arbitrary public low-degree
polynomial to $L$ breaks verification unless the verifier has a matching
authenticated mask evaluation. The required conditional interface is described in \cref{sec:mask}, but
its implementation from the forests is not specified or proved here.
\item Decide whether hiding is needed for more than one opening. If yes,
fresh independent commitments alone are not a repair unless their witnesses
are linked by a new proof; reusing punctured forests discloses their witness.
\item Analyze the concrete puncturing counterexample and specify any proposed
replacement consistency mechanism. Committing additional pre-tree data may
help, but requires new opening, leakage, and extraction analysis.
\item Confirm the coordinatewise lift, and choose explicitly between a
proved valid-encoding extraction guarantee and the encoding-checked mode.
For non-power-of-two $n$, range validity cannot be omitted without justification.
\item Prove a commitment-time malicious-consistency/extraction statement
before applying the conditional polynomial bound to arbitrary commitments.
\item Use the revised stacking definitions to identify whether the intended
compiler needs linked witnesses across distinct commitments. Specify
its exact transcript before claiming a composition result.
\end{enumerate}

\paragraph{Suggested eventual research-state updates; not applied.}
Record (i) the coefficient-testing identity and four-entry separator with
its concrete constant-advantage bound; (ii) the two-opening witness disclosure;
(iii) the distinction between packed blocks and coordinatewise extension-field
evaluations; (iv) the encoding/range constraints and their degree cost;
(v) the span-rank distribution of $\Delta$ and the conditional bound
\eqref{eq:bound}; and (vi) the adaptive pre-tree-key counterexample and the resulting failure
of structural common-witness consistency. Literature proof transfer
does not establish the remaining cryptographic claims. The original conjectures about binding
and online extraction have not become established results.

\begin{center}
\begin{tabular}{p{0.49\linewidth}p{0.41\linewidth}}
\toprule
Conclusion & Status\\
\midrule
Homogenization and affine correctness identities & Established algebra\\
Coefficient-based candidate testing & Established algebra\\
Index distinction when $\beta v_1\ne0$ & Established conditional separator\\
Single-opening index hiding & Broken with constant advantage\\
Two-opening witness recovery & Established for honest forests\\
Polynomial soundness for a fixed affine witness & Conditional bound proved here\\
Pre-tree enforcement of common forest sums & Refuted by explicit adaptive keys\\
Classical commitment-time online extraction & Unresolved\\
QROM extraction and Fiat--Shamir security & Unresolved\\
Concrete forest algorithms / revised stacking comparison & Specified / technical sources inspected\\
Full index hiding / complete secure stacking & Candidate insecure / compiler unresolved\\
\bottomrule
\end{tabular}
\end{center}

\section{Validation and remaining proof obligations}\label{sec:validation}
The manuscript uses explicit least-significant-bit folding and most-significant
root-splitting conventions, with the pre-challenge in \cref{eq:prechallenge}
connecting them. Field operations in polynomial checking are coordinatewise;
XOR applies only to binary node strings or binary digest strings.
The honest reconstruction proof includes $D=1$ and $\tau=1$.
The security conclusions separate forest acceptance from full opening
acceptance, and algebraic identities from cryptographic extraction claims.

Deterministic computational checks used 8-bit node values and the same
XOR-preserving rule (with domain-separated SHA-256 output truncation):
398 complete challenge layouts across $(\tau,D)=(1,1),(1,3),(2,1),
(2,3),(3,2),(4,2)$ passed both honest reconstruction and the adaptive
pre-tree-key counterexample, including missing-digest reconstruction.
In $\mathbb F_{16}$ with modulus $U^4+U+1$, 2,000 random valid-witness
instances with $q=1,\ldots,4$ verified the formal leading coefficient
and all 16 evaluation points. Exhaustive three-mask-coordinate checks
verified the checked-mode separator for sampled challenge coefficients.
The span-mass identity was checked exhaustively for $M=1,\ldots,4$
over $\mathbb F_4$; the excluded-index example was checked on all eight
three-entry binary opening vectors.
The source was built with pdfLaTeX and Biber; references were resolved
and the resulting PDF pages were rendered for visual inspection.
These checks supplement the proofs; they do not establish cryptographic
security for unbounded parameters or replace missing ROM/QROM reductions.
The coefficient attack already contradicts the requested hiding target.
The immediate research priority is a masking design with an authenticated
mask evaluation, followed by a malicious-forest analysis that does not assume
an uncommitted pre-tree enforces common sums. Neither modification has been
adopted as a repaired construction in this manuscript.

\printbibliography
\end{document}
